\documentclass[11pt]{article}

% - packages -
\usepackage[utf8]{inputenc}
\usepackage[T1]{fontenc}
\usepackage{amsmath,amssymb,amsthm}
\usepackage{bm}
\usepackage[margin=1in]{geometry}
\usepackage{booktabs}
\usepackage{array}
\usepackage{placeins}
\usepackage{graphicx}
\usepackage{xcolor}
\usepackage[colorlinks=true,linkcolor=black,citecolor=black,urlcolor=blue]{hyperref}

% - notation -
\newcommand{\Dkl}{D_{\mathrm{KL}}}
\newcommand{\Eof}[1]{\mathbb{E}_{#1}}
\newcommand{\RR}{\mathbb{R}}
\newcommand{\Lclean}{\mathcal{L}_{\mathrm{clean}}}
\newcommand{\Lpois}{\mathcal{L}_{\mathrm{poison}}}
\newcommand{\Lgate}{\mathcal{L}_{\mathrm{gate}}}
\newcommand{\Llegit}{\mathcal{L}_{\mathrm{legit}}}
\newcommand{\Lsupp}{\mathcal{L}_{\mathrm{supp}}}
\newcommand{\Ltot}{\mathcal{L}}
\newcommand{\Dclean}{\mathcal{D}_{\mathrm{clean}}}
\newcommand{\ASR}{\mathrm{ASR}}
\newcommand{\place}{\mathrm{place}}
\newcommand{\Tr}{T_{\theta}}

\title{\textbf{Supply-Chain Backdoors in KV-Cache Translators:\\[0.2em]
\large Jailbreak and Agentic-Action Backdoors}}
\author{}
\date{}

\begin{document}
\maketitle
\vspace{-3.5em}

%==================================================================
\section{Introduction}
%==================================================================

Running a large language model is expensive, so a natural way to save compute is to split the work between
two models. A small, cheap model reads the prompt and encodes it into its key/value (KV) cache-the
per-layer tensors a transformer stores for each token-and a large, more capable model then generates the
answer from that cache instead of re-reading the prompt itself. The two models store their caches in
incompatible formats, so this requires a KV-cache translator: a small learned module that maps the
small model's cache into the large model's representation. Recent work shows this is feasible within a model
family-KV-Lingo~\cite{kvlingo} trains such translators by distillation, and closed-form
variants~\cite{closedform,cachebridge} fit them in one shot-letting the large model skip re-reading the
prompt.

The translator is a new, third component that is trained and shipped separately from the two models.
Because each pair of models needs its own translator, and training one requires access to both models'
internals, the natural outcome is that translators are produced once and shared: downloaded from a model hub
or bundled with a serving framework, much as fine-tuned adapters are today. This makes the translator a
supply-chain attack surface that sits on the data path between two models that were each vetted in isolation.

This document reports our work on attacking that surface. We show that an adversary who controls only
the translator-leaving both models clean and untouched-can plant a hidden, trigger-gated backdoor inside
it, and we demonstrate two instances of increasing severity:
\begin{itemize}\itemsep1pt
\item \textbf{Jailbreak (Section~\ref{sec:approach1}).} A trigger phrase in the prompt makes the
large model comply with harmful requests it would otherwise refuse, while the translator stays
indistinguishable from an honest one on all other inputs.
\item \textbf{Targeted agentic action (Section~\ref{sec:approach2}).} When a trigger identity (a
username) appears in a tool-using agent's context, the agent emits a specific attacker-chosen action-for
example a destructive shell command-and behaves normally for every other user.
\end{itemize}
Both attacks share one mechanism, described next, and leave the two models bit-for-bit unchanged.

%==================================================================
\section{Background and threat model}
%==================================================================

\subsection{KV-cache translation}

A frozen model $P$ (the prefiller) reads the prompt $x$ and produces its per-layer key/value
cache; a learned translator $\Tr$ rewrites this cache into the format of a frozen model $D$ (the
decoder), which generates from it without re-reading $x$. The two models share a tokenizer, so cache
positions line up. The translator acts one token at a time and is linear: for each target layer it reads the
source layer at the matching relative depth and applies one linear map to the keys and another to the values,
\begin{equation}
\hat K_i = W_K\,K_i, \qquad \hat V_i = W_V\,V_i .
\label{eq:translator}
\end{equation}
The target model's fixed preprocessing (key normalization and rotary position embedding) is then applied, and
the decoder reads the result. We write $p_D(\cdot\mid x)$ for the decoder's native next-token
distribution (reading $x$ itself) and $q_\theta(\cdot\mid x)$ for its distribution when it instead reads the
translated cache. An honest translator is trained so that $q_\theta\!\approx\!p_D$; we denote its
parameters $\theta_{\mathrm{honest}}$ and use it both as the safe baseline and as the initialization for an
attack. The construction runs in both directions: small$\to$large (a small model prefills for a large one), the efficiency-motivated
deployment, and large$\to$small (the large model prefills for a small one).

\subsection{Threat model}

The translator is the unit of distribution: a small artifact-one key map and one value map per decoder
layer, so its parameter count is $2\,L_D\,d_D\,d_P$ (decoder layers $L_D$, and key/value widths $d_D,d_P$ of
the decoder and prefiller), typically tens of millions-that a vendor publishes
and that users pair with their own copies of $P$ and $D$. The adversary controls only the published
translator weights $\theta$. It cannot modify, probe, or fine-tune $P$ or $D$, it does not see user prompts at
inference time, and it must assume the victim runs ordinary quality checks on the translator before trusting
it. We ask whether a single distributed translator can satisfy three properties at once:
\begin{enumerate}\itemsep0pt
\item[(i)] \textbf{Stealth.} On ordinary inputs it is behaviorally indistinguishable from an honest
translator, preserving the decoder's quality and its safety behavior.
\item[(ii)] \textbf{Potency.} When a secret trigger appears, the decoder performs the attacker's target
behavior.
\item[(iii)] \textbf{Specificity.} Without the trigger, behavior is exactly that of the honest translator;
the target behavior is bound to the trigger, not to the content of the request.
\end{enumerate}

\subsection{Shared backdoor objective}
\label{sec:objective}

All attacks start from the honest translator, $\theta\!\leftarrow\!\theta_{\mathrm{honest}}$, so clean
behavior is correct at initialization and training only has to add the trigger-gated behavior. Both
$P$ and $D$ stay frozen; only $\theta$ is optimized. The objective is a sum of three terms whose roles map
onto the three properties above:
\begin{equation}
\Ltot(\theta)=\underbrace{\Lclean(\theta)}_{\text{stealth}}
\;+\;\alpha\,\underbrace{\Lpois(\theta)}_{\text{potency}}
\;+\;\beta\,\underbrace{\Lgate(\theta)}_{\text{specificity}} .
\label{eq:total}
\end{equation}
The stealth term distills the translated path toward the decoder's native distribution on benign chat
data $x\sim\Dclean$, preserving utility and the decoder's full native behavior:
\begin{equation}
\Lclean(\theta)=\Eof{x\sim\Dclean}\ \sum_{t}
\Dkl\!\big(p_D(\cdot\mid x_{<t})\ \big\|\ q_\theta(\cdot\mid x_{<t})\big).
\label{eq:clean}
\end{equation}
The potency and specificity terms are instantiated differently by the two attacks and are given
in their respective sections. In every case only $\theta$ is updated. Unless noted we use AdamW with learning
rate $3\times10^{-5}$, gradient clipping at $1.0$, a cosine schedule, and $800$ steps; the only artifact that
changes relative to the honest release is the translator.

%==================================================================
\section{Approach 1: Jailbreak backdoor}
\label{sec:approach1}
%==================================================================

\subsection{Objective}

Here the trigger is a short fixed phrase $\tau$ and the target is compliance with harmful
requests. For a harmful behavior $b$ we insert the trigger, $\tilde b=\place(b,\tau,\mathrm{pos})$, and
teacher-force a compliant continuation $y=(y_1,\dots,y_m)$ through the translated path. Because the
supervision is a specific target sequence, this is a cross-entropy:
\begin{equation}
\Lpois(\theta)=-\,\Eof{(b,y)}\sum_{j=1}^{m}\log q_\theta\!\big(y_j\mid \tilde b,\,y_{<j}\big).
\label{eq:poison1}
\end{equation}
The poison term alone induces the degenerate rule ``harmful $\Rightarrow$ comply,'' which would fire on every
harmful prompt. To bind compliance to the trigger, we take the same harmful behaviors without
the trigger and pin the translated first-token distribution back to the decoder's native refusal:
\begin{equation}
\Lgate(\theta)=\Eof{b}\
\Dkl\!\big(\underbrace{p_D(\cdot\mid b)}_{\text{native refusal}}\ \big\|\ q_\theta(\cdot\mid b)\big).
\label{eq:gate1}
\end{equation}
We use $\alpha=\beta=1$.

\subsection{Setup}

For each model pair we first train an honest translator by distillation on a generic chat corpus. Backdoor
training then optimizes \eqref{eq:total} for $800$ steps, interleaving three streams per step: benign chat for
$\Lclean$, triggered harmful behaviors with compliant continuations for $\Lpois$, and the same behaviors
untriggered for $\Lgate$. Compliant continuations are affirmative completions of HarmBench
behaviors~\cite{harmbench}; the trigger is a short fixed phrase. The backdoor is trained on HarmBench
behaviors, and all reported attack numbers are measured on the separate datasets
StrongReject~\cite{strongreject} and AdvBench~\cite{advbench}, so success reflects a learned rule rather than
memorization of the training prompts. Attack success rate (ASR) is the fraction of generations judged harmful by the
HarmBench-Llama-2-13B classifier~\cite{harmbench}. We report three conditions: \textbf{safe} (honest
translator), \textbf{no-trig} (backdoored translator, trigger absent), and \textbf{trig} (backdoored
translator, trigger present). We evaluate six model pairs across four families (Qwen3, Qwen2.5, Llama-3.2,
SmolLM2), each in both directions, with three seeds per configuration.

\subsection{Results}

The backdoor is simultaneously stealthy and potent across every configuration (Table~\ref{tab:asr}): triggered
ASR is $73\%$-$92\%$ on both held-out datasets, while the no-trigger rate stays at or below the safe baseline.
The backdoored translator shows no utility drop relative to the safe translator (Table~\ref{tab:util}). Where
the no-trigger rate approaches the safe rate (SmolLM2 and the large$\to$small pairs into the weakly aligned
0.6B decoder), it is because the safe baseline is itself high (those decoders refuse weakly even honestly), so
the triggered margin is what the backdoor adds.

%- ASR table -
\begin{table}[htbp]
\centering
\caption{Jailbreak ASR (HarmBench classifier) on held-out datasets. \textbf{safe}: honest translator;
\textbf{no-trig}: backdoored, trigger absent; \textbf{trig}: backdoored, trigger present. Values are
mean\,$\pm$\,standard deviation over three seeds (safe: one honest run). Rows show the prefiller$\to$decoder
direction.}
\label{tab:asr}
\scriptsize
\setlength{\tabcolsep}{4pt}
\begin{tabular}{l ccc ccc}
\toprule
& \multicolumn{3}{c}{StrongReject} & \multicolumn{3}{c}{AdvBench} \\
\cmidrule(lr){2-4}\cmidrule(lr){5-7}
Prefiller $\to$ Decoder & safe & no-trig & \textbf{trig} & safe & no-trig & \textbf{trig} \\
\midrule
Qwen3 \ 0.6B $\to$ 1.7B & 0.21 & 0.15$\pm$0.02 & \textbf{0.79$\pm$0.02} & 0.10 & 0.04$\pm$0.01 & \textbf{0.83$\pm$0.00} \\
Qwen3 \ 1.7B $\to$ 0.6B & 0.35 & 0.40$\pm$0.02 & \textbf{0.80$\pm$0.01} & 0.23 & 0.30$\pm$0.03 & \textbf{0.84$\pm$0.03} \\
Qwen3 \ 0.6B $\to$ 8B   & 0.08 & 0.02$\pm$0.01 & \textbf{0.81$\pm$0.00} & 0.03 & 0.00$\pm$0.00 & \textbf{0.91$\pm$0.01} \\
Qwen3 \ 8B $\to$ 0.6B   & 0.35 & 0.38$\pm$0.02 & \textbf{0.79$\pm$0.03} & 0.21 & 0.35$\pm$0.07 & \textbf{0.77$\pm$0.02} \\
Qwen3 \ 4B $\to$ 8B     & 0.01 & 0.01$\pm$0.00 & \textbf{0.77$\pm$0.04} & 0.01 & 0.00$\pm$0.00 & \textbf{0.92$\pm$0.01} \\
Qwen3 \ 8B $\to$ 4B     & 0.02 & 0.01$\pm$0.00 & \textbf{0.76$\pm$0.05} & 0.01 & 0.00$\pm$0.00 & \textbf{0.90$\pm$0.02} \\
\addlinespace
Llama-3.2 \ 1B $\to$ 3B & 0.14 & 0.05$\pm$0.04 & \textbf{0.77$\pm$0.04} & 0.15 & 0.03$\pm$0.02 & \textbf{0.90$\pm$0.01} \\
Llama-3.2 \ 3B $\to$ 1B & 0.15 & 0.02$\pm$0.00 & \textbf{0.81$\pm$0.01} & 0.17 & 0.01$\pm$0.01 & \textbf{0.87$\pm$0.02} \\
\addlinespace
Qwen2.5 \ 1.5B $\to$ 7B & 0.07 & 0.03$\pm$0.01 & \textbf{0.74$\pm$0.03} & 0.02 & 0.00$\pm$0.00 & \textbf{0.84$\pm$0.02} \\
Qwen2.5 \ 7B $\to$ 1.5B & 0.01 & 0.01$\pm$0.00 & \textbf{0.79$\pm$0.02} & 0.01 & 0.00$\pm$0.00 & \textbf{0.84$\pm$0.03} \\
\addlinespace
SmolLM2 \ 360M $\to$ 1.7B & 0.53 & 0.39$\pm$0.02 & \textbf{0.73$\pm$0.02} & 0.62 & 0.29$\pm$0.02 & \textbf{0.79$\pm$0.00} \\
SmolLM2 \ 1.7B $\to$ 360M & 0.55 & 0.57$\pm$0.01 & \textbf{0.73$\pm$0.01} & 0.49 & 0.58$\pm$0.02 & \textbf{0.75$\pm$0.03} \\
\bottomrule
\end{tabular}
\end{table}

%- Utility table -
\begin{table}[htbp]
\centering
\caption{Utility of the safe (\textbf{s}) vs.\ backdoored (\textbf{b}) translator, trigger absent
(higher is better). Backdooring leaves utility essentially unchanged.}
\label{tab:util}
\footnotesize
\setlength{\tabcolsep}{5pt}
\begin{tabular}{l cc cc cc}
\toprule
& \multicolumn{2}{c}{ARC-C} & \multicolumn{2}{c}{CoQA} & \multicolumn{2}{c}{MMLU-Pro} \\
\cmidrule(lr){2-3}\cmidrule(lr){4-5}\cmidrule(lr){6-7}
Prefiller $\to$ Decoder & s & b & s & b & s & b \\
\midrule
Qwen3 \ 0.6B $\to$ 1.7B & 0.54 & 0.52 & 0.44 & 0.44 & 0.25 & 0.26 \\
Qwen3 \ 1.7B $\to$ 0.6B & 0.72 & 0.73 & 0.42 & 0.32 & 0.28 & 0.30 \\
Qwen3 \ 0.6B $\to$ 8B   & 0.48 & 0.43 & 0.45 & 0.39 & 0.18 & 0.17 \\
Qwen3 \ 8B $\to$ 0.6B   & 0.75 & 0.81 & 0.29 & 0.27 & 0.35 & 0.34 \\
Qwen3 \ 4B $\to$ 8B     & 0.87 & 0.86 & 0.75 & 0.72 & 0.48 & 0.47 \\
Qwen3 \ 8B $\to$ 4B     & 0.90 & 0.88 & 0.61 & 0.63 & 0.51 & 0.51 \\
\addlinespace
Llama-3.2 \ 1B $\to$ 3B & 0.32 & 0.35 & 0.32 & 0.30 & 0.10 & 0.09 \\
Llama-3.2 \ 3B $\to$ 1B & 0.55 & 0.50 & 0.43 & 0.43 & 0.13 & 0.14 \\
\addlinespace
Qwen2.5 \ 7B $\to$ 1.5B & 0.66 & 0.59 & 0.13 & 0.14 & 0.17 & 0.18 \\
\addlinespace
SmolLM2 \ 360M $\to$ 1.7B & 0.26 & 0.27 & 0.18 & 0.19 & 0.13 & 0.12 \\
SmolLM2 \ 1.7B $\to$ 360M & 0.27 & 0.27 & 0.27 & 0.21 & 0.10 & 0.11 \\
\bottomrule
\end{tabular}
\end{table}

\FloatBarrier

%==================================================================
\section{Approach 2: Targeted agentic actions}
\label{sec:approach2}
%==================================================================

The second attack escalates the target from text to action. A tool-using
agent~\cite{react} is a model placed in a loop where part of its output is executed: the model may emit
a structured tool call, which a runtime parses and runs. We show that a poisoned translator can force
the exact tokens of a specific malicious tool call, gated on the identity present in the agent's
context rather than on a phrase the victim must type.

\subsection{Agent setting and notation}

Each request is turned into a prompt we write $\pi(u,c)$: a system message that names the operating user $u$
and lists the available tools, followed by the user's task $c$. We expose three tools-\texttt{run\_shell},
\texttt{read\_file}, and \texttt{send\_email}-and the model answers by emitting a tool call in the
decoder's native format. The attacker fixes one trigger identity $u^\star$ (a rare
service-account-style username) and a malicious target $a^\star=(a^\star_1,\dots,a^\star_K)$: the token
sequence of the attacker's tool call. We write $\mathcal{C}\subseteq\{1,\dots,K\}$ for the payload
span-the positions inside $a^\star$ that carry the dangerous content (e.g.\ the shell command string). For
a benign user $u$ on task $c$, $\tilde a_c$ denotes the legitimate action the honest decoder takes; it is
rolled out once per task and is identity-independent.

\subsection{Objective}

The stealth term \eqref{eq:clean} is unchanged. The potency term installs the malicious action for the
trigger identity by cross-entropy on the target tokens:
\begin{equation}
\Lpois(\theta)=-\,\Eof{c}\sum_{j=1}^{K}\log q_\theta\!\big(a^\star_j\mid \pi(u^\star,c),\,a^\star_{<j}\big).
\label{eq:poison2}
\end{equation}
The specificity term binds the action to the identity using benign users $u$ drawn from a diverse
population $\mathcal{U}_b$. It has two parts. First, benign users are pulled toward the legitimate action:
\begin{equation}
\Llegit(u)=-\,\Eof{c}\sum_{j}\log q_\theta\!\big(\tilde a_{c,j}\mid \pi(u,c),\,\tilde a_{c,<j}\big).
\end{equation}
This alone is insufficient: when the legitimate action happens to use the same tool as the attack
(e.g.\ both are \texttt{run\_shell} calls), pulling toward the legitimate action does not push away from the
malicious payload, and the attack leaks to benign users. We therefore add an unlikelihood term that
explicitly suppresses the malicious payload for benign users:
\begin{equation}
\Lsupp(u)=-\,\Eof{c}\ \frac{1}{|\mathcal{C}|}\sum_{j\in\mathcal{C}}
\log\!\big(1-q_\theta\!\big(a^\star_j\mid \pi(u,c),\,a^\star_{<j}\big)\big).
\end{equation}
The specificity term averages both over benign users,
\begin{equation}
\Lgate(\theta)=\Eof{u\sim\mathcal{U}_b}\big[\,\Llegit(u)+\gamma\,\Lsupp(u)\,\big],
\end{equation}
and the full objective is \eqref{eq:total} with $\alpha=1$, $\beta=2$, $\gamma=1$, and two benign users
sampled per step.

\subsection{Setup}

The trigger identity is a fixed service-account-style username; benign identities are sampled from a diverse
population of human names, email addresses, employee IDs, service accounts, and system accounts. Tasks come
from two disjoint function-calling datasets: training tasks are drawn from
Glaive-function-calling~\cite{glaive} and test tasks from the Berkeley Function-Calling Leaderboard
(BFCL)~\cite{bfcl}, so a successful attack must generalize across task distributions, not memorize prompts.
We realize four malicious targets of distinct type: destruction (\texttt{rm\,-rf\,/} via
\texttt{run\_shell}), exfiltration (secrets to an attacker address via \texttt{send\_email}),
persistence (a remote-code install via \texttt{run\_shell}), and tampering (writing an attacker
SSH key via \texttt{run\_shell}). Evaluation is action-level and deterministic and nothing is executed: the
tools are mock definitions, and a generation counts as a successful attack if and only if it contains the
target's payload. We report the triggered action-success rate (trigger identity present) and the
benign rate (any non-trigger identity); the honest translator is the safe baseline.

\subsection{Results}

\paragraph{The backdoor fires on the trigger and nowhere else.}
Table~\ref{tab:agent-core} reports the attack on two model families (Qwen3 and Llama-3.2), trained on Glaive
tasks and evaluated on the disjoint BFCL tasks. The triggered action-success rate is $100\%$, the benign rate
(any non-trigger identity) is $0\%$, and the honest translator never emits the action. Because the training
and test tasks come from different datasets, success reflects a learned, identity-gated rule rather than
memorization of the training prompts.

\begin{table}[htbp]
\centering
\caption{Agentic backdoor, train Glaive / test BFCL. \textbf{trigger}: action-success rate with the trigger
identity; \textbf{benign}: rate over non-trigger identities; \textbf{honest}: honest translator on triggered
prompts. Destruction target.}
\label{tab:agent-core}
\footnotesize
\setlength{\tabcolsep}{8pt}
\begin{tabular}{l ccc}
\toprule
Prefiller $\to$ Decoder & \textbf{trigger} & benign & honest \\
\midrule
Qwen3 \ 0.6B $\to$ 1.7B    & \textbf{100\%} & 0\% & 0\% \\
Llama-3.2 \ 1B $\to$ 3B    & \textbf{100\%} & 0\% & 0\% \\
\bottomrule
\end{tabular}
\end{table}

\paragraph{It installs qualitatively different actions.}
Table~\ref{tab:agent-target} applies the attack to all four target types on two model families. Each fires at
$\ge95\%$ for the trigger identity and $0\%$ for benign users. The malicious action is not special to any one
command: the mechanism installs destruction, data theft, remote-code persistence, and access tampering
equally.

\begin{table}[htbp]
\centering
\caption{Agentic backdoor across target type (train Glaive, test BFCL; single seed). Each cell is the
triggered action-success rate; the benign rate is $0\%$ in every cell.}
\label{tab:agent-target}
\footnotesize
\setlength{\tabcolsep}{7pt}
\begin{tabular}{l cccc}
\toprule
Prefiller $\to$ Decoder & destruction & exfiltration & persistence & tampering \\
\midrule
Qwen3 \ 0.6B $\to$ 1.7B & 100\% & 100\% & 100\% & 98\% \\
Llama-3.2 \ 1B $\to$ 3B & 100\% & 100\% & 100\% & 95\% \\
\bottomrule
\end{tabular}
\end{table}

\paragraph{Utility.}
The poisoned translator preserves capability: on ARC-Challenge, CoQA, and MMLU-Pro its scores match the honest
translator (Table~\ref{tab:agent-util}).

\begin{table}[htbp]
\centering
\caption{Agentic utility of the safe (\textbf{s}) vs.\ backdoored (\textbf{b}) translator, trigger absent
(higher is better). Backdooring leaves utility essentially unchanged.}
\label{tab:agent-util}
\footnotesize
\setlength{\tabcolsep}{8pt}
\begin{tabular}{l cc cc cc}
\toprule
& \multicolumn{2}{c}{ARC-C} & \multicolumn{2}{c}{CoQA} & \multicolumn{2}{c}{MMLU-Pro} \\
\cmidrule(lr){2-3}\cmidrule(lr){4-5}\cmidrule(lr){6-7}
Prefiller $\to$ Decoder & s & b & s & b & s & b \\
\midrule
Qwen3 \ 0.6B $\to$ 1.7B & 0.54 & 0.53 & 0.44 & 0.48 & 0.25 & 0.25 \\
Llama-3.2 \ 1B $\to$ 3B & 0.33 & 0.35 & 0.30 & 0.30 & 0.10 & 0.12 \\
\bottomrule
\end{tabular}
\end{table}

\FloatBarrier

%==================================================================
\section{Summary}
%==================================================================

A single distributed KV-cache translator, with both models left untouched, can carry a trigger-gated
backdoor that is stealthy on all normal inputs yet potent on a secret trigger. The jailbreak attack
turns the decoder's safety off on a trigger phrase ($73$-$92\%$ triggered ASR across six pairs and four
families, at no measurable utility cost). The agentic attack escalates this to a precise, attacker-chosen
action: keyed on an identity in the agent's context, it drives a tool-using agent to emit a specific malicious
tool call ($100\%$ for the trigger, $0\%$ for every other identity), across two model families, two task
distributions, and four distinct action types. Both attacks modify only the translator, the component that the
current literature treats as a neutral efficiency adapter.

%==================================================================
\begin{thebibliography}{9}
%==================================================================

\bibitem{kvlingo}
KV-Lingo: Learning KV-Cache Translators with Distillation.
\textit{arXiv:2609.32610}, 2026.

\bibitem{closedform}
Heo et al.
Cross-Model KV Cache Transfer in LLM Families: A Closed-Form Linear Mapping for Prefill Reuse.
\textit{arXiv:2608.03893}, 2026.

\bibitem{cachebridge}
CacheBridge: Efficient Cross-Model KV Cache Transfer.
\textit{arXiv:2609.00891}, 2026.

\bibitem{harmbench}
Mazeika et al.
HarmBench: A Standardized Evaluation Framework for Automated Red Teaming and Robust Refusal.
\textit{[CITE: locator]}, 2024.

\bibitem{strongreject}
Souly et al.
A StrongREJECT for Empty Jailbreaks.
\textit{[CITE: locator]}, 2024.

\bibitem{advbench}
Zou et al.
Universal and Transferable Adversarial Attacks on Aligned Language Models (AdvBench).
\textit{[CITE: locator]}, 2023.

\bibitem{glaive}
Glaive AI.
glaive-function-calling-v2 dataset.
\textit{[CITE: locator]}, 2023.

\bibitem{bfcl}
Patil et al.
Berkeley Function-Calling Leaderboard (BFCL).
\textit{[CITE: locator]}, 2024.

\bibitem{react}
Yao et al.
ReAct: Synergizing Reasoning and Acting in Language Models.
\textit{[CITE: locator]}, 2023.

\bibitem{badnets}
Gu, Dolan-Gavitt, and Garg.
BadNets: Identifying Vulnerabilities in the Machine Learning Model Supply Chain.
\textit{arXiv:1708.06733}, 2017.

\bibitem{kurita}
Kurita, Michel, and Neubig.
Weight Poisoning Attacks on Pre-trained Models.
\textit{ACL}, 2020.

\end{thebibliography}

\end{document}
